\subsection{Criterios de Nagata para que $\operatorname{Reg}(X)$ sea abierto.}
\label{IV.2.s6.12-fr}

\begin{env}[6.12.1]
\label{IV.2.s6.12.1-fr}
Dado un preesquema localmente noetheriano $X$, se llama \emph{lugar
singular} de $X$ y se denota $\operatorname{Sing}(X)$ el conjunto de los puntos
$x\in X$ en los que $X$ no es regular en el punto $x$ (es decir,
aquellos para los que el anillo local $\mathcal O_x$ no es regular); el
complemento $X\setmin\operatorname{Sing}(X)$, conjunto de los $x\in X$ en los que $X$
es regular, se denota $\operatorname{Reg}(X)$. En este apartado nos proponemos
dar condiciones bajo las cuales $\operatorname{Sing}(X)$ es \emph{cerrado} (es decir,
$\operatorname{Reg}(X)$ es \emph{abierto}).
\end{env}

\begin{proposition}[6.12.2]
\label{IV.2.s6.12.2-fr}
Sea $X$ un preesquema localmente noetheriano. Las condiciones siguientes son
equivalentes:
\begin{enumerate}
\item[\textup{a)}] $\operatorname{Reg}(X)$ es abierto en $X$.
\item[\textup{b)}] Para todo $x\in\operatorname{Reg}(X)$, existe una parte
abierta no vacía de $\overline{\{x\}}$ contenida en
$\operatorname{Reg}(X)$.
\end{enumerate}
Además, estas condiciones se siguen de la siguiente:
\begin{enumerate}
\item[\textup{c)}] Para todo $x\in\operatorname{Reg}(X)$, si se designa por $Y$ el
subpreesquema cerrado reducido de $X$ cuyo espacio subyacente es
$\overline{\{x\}}$, entonces $\operatorname{Reg}(Y)$ es una vecindad de $x$ en $Y$.
\end{enumerate}
\end{proposition}

\begin{proof}
La equivalencia de \emph{a)} y \emph{b)} se deduce de que $\operatorname{Reg}(X)$ es
estable por generización (\textbf 0, 17.3.2) y de (\textbf{0}$_{\rm III}$, 9.2.6). Para
probar que \emph{c)} implica \emph{b)}, es posible limitarse al caso en que
$X=\operatorname{Spec}(A)$ es una vecindad abierta afín de $x$; si
$(t_i)_{1\leq i\leq n}$ es un sistema regular de parámetros (\textbf 0, 17.1.6)
del anillo local regular $A_x$, se puede suponer (reemplazando, si es
necesario, $X$ por una vecindad abierta de $x$) que $t_i=(s_i)_x$ para todo
$i$, donde $s_i\in A$ y la familia $(s_i)_{1\leq i\leq n}$ es $A$-regular
(\textbf 0, 15.2.4). Entonces $Y=\operatorname{Spec}(A/\mathfrak p)$, donde
$\mathfrak p=\mathfrak j_x$; como los $t_i$ generan el ideal maximal
$\mathfrak m=\mathfrak pA_x$ de $A_x$, se puede suponer también (reemplazando
$X$ por una vecindad más pequeña de $x$) que los $s_i$ generan
$\mathfrak p$. Para todo $y\in Y$, los $(s_i)_y$ generan, por tanto,
$\mathfrak p_y$; como forman una sucesión $A_y$-regular, forman
parte de un sistema de parámetros de $A_y$ (\textbf 0, 16.4.1); por consiguiente,
se deduce de (\textbf 0, 17.1.7) que, si $(A/\mathfrak p)_y=A_y/\mathfrak p_y$
es regular, también lo es $A_y$, de donde se sigue la conclusión.
\end{proof}
